\documentclass[11pt, a4paper]{article}
\usepackage{fontspec}
\usepackage[turkish]{babel}
\setmainfont{Latin Modern Roman}
\usepackage[margin=2.5cm]{geometry}
\usepackage{hyperref}
\usepackage[shortlabels]{enumitem}
\usepackage{xcolor}

% Link ve PDF ayarları
\hypersetup{
    colorlinks=true,
    linkcolor=blue,
    filecolor=magenta,      
    urlcolor=cyan,
    pdftitle={System Requirements Document},
    pdfpagemode=FullScreen,
}

% Kapak Sayfası Bilgileri
\title{
    \vspace{2cm}
    \Huge \textbf{System Requirements Document}\\
    \vspace{1cm}
    \LARGE \textbf{Aegis Command \& Control}\\
    \vspace{2cm}
}
\author{\Large \textbf{Sürüm:} 1.0.0}
\date{\vspace{1cm} \Large \today}

\begin{document}

\maketitle
\thispagestyle{empty}

\vspace{3cm}
\begin{center}
    \Large \textbf{Target Hardware:} STM32F407G-DISC1 \\
    \vspace{0.5cm}
    \Large \textbf{Software Stack:} C++20, FreeRTOS, Node.js, Next.js
\end{center}

\newpage
\pagenumbering{roman}
\tableofcontents
\newpage
\pagenumbering{arabic}

\section{Introduction}

\subsection{Purpose}
Bu doküman, STM32 mikrodenetleyicisi ve web tabanlı bir komuta kontrol arayüzünden oluşan "Aegis Command \& Control" sisteminin donanım, yazılım ve haberleşme gereksinimlerini tanımlar.

\subsection{System Overview}
AC2 Sistemi, üç ana katmandan oluşan bir "Full-Stack Embedded" mimarisidir:
\begin{enumerate}[1.]
    \item \textbf{Edge Node (STM32):} Deterministik durum makinesini (State Machine) işleten, sensör/donanım etkileşimini sağlayan gerçek zamanlı çekirdek.
    \item \textbf{Gateway Bridge:} Seri port (UART) üzerinden gelen paketleri çözümleyip WebSockets (TCP/IP) üzerinden arayüze aktaran ara katman.
    \item \textbf{Command Control Interface (Next.js):} Kullanıcının sistemi uzaktan izlediği ve kontrol ettiği, yüksek tepkime hızına sahip modern web arayüzü.
\end{enumerate}

\section{System Architecture \& Constraints}

\noindent \textbf{ARCH-01: SOLID \& Design Patterns}
\begin{itemize}
    \item Gömülü sistem yazılımı; Dependency Inversion (DIP) ve Single Responsibility (SRP) prensiplerine kesin olarak uyacaktır.
    \item Durum yönetimi için \textit{State Pattern}, kesme yönetimi için \textit{Observer Pattern}, donanım soyutlaması için \textit{Wrapper/Adapter Pattern} kullanılacaktır.
\end{itemize}

\vspace{0.5cm}
\noindent \textbf{ARCH-02: Memory \& Resource Management}
\begin{itemize}
    \item STM32 üzerinde çalışma zamanında (runtime) dinamik bellek tahsisi (\texttt{new}, \texttt{malloc}) kesinlikle yapılmayacaktır. 
    \item Tüm nesneler statik olarak tahsis edilecek, bellek sızıntısı ve fragmentasyon riskleri donanım seviyesinde engellenecektir.
\end{itemize}

\vspace{0.5cm}
\noindent \textbf{ARCH-03: RTOS Integration}
\begin{itemize}
    \item Görevler (Tasks) arası iletişim, paylaşılan bellek yerine FreeRTOS Queue ve Semaphore yapıları kullanılarak Thread-Safe şekilde yapılacaktır.
\end{itemize}

\section{Functional Requirements (FR)}

\subsection{Embedded Core (STM32) Requirements}

\noindent \textbf{FR-01: State Management}\\
Sistem üç ana durumu (IDLE, SEARCH, TRACK) yönetecek ve her durum değişimi Gateway'e anlık olarak raporlanacaktır.
\begin{itemize}
    \item \textbf{IDLE:} Sistem beklemededir, LED'ler kapalıdır. Dışarıdan "START" komutu bekler.
    \item \textbf{SEARCH:} Sistem arama modundadır. Yeşil LED (PD12) 500ms periyotla yanıp söner. Bir hedef bulunduğunda TRACK moduna geçer.
    \item \textbf{TRACK:} Sistem kilitlenmiştir. Kırmızı LED (PD14) sürekli yanar. Hedef kaybedilirse SEARCH moduna, "STOP" komutu gelirse IDLE moduna döner.
\end{itemize}

\vspace{0.5cm}
\noindent \textbf{FR-02: Hardware Interrupts}
\begin{itemize}
    \item Kullanıcı butonu (User Button), sisteme donanımsal kesme (EXTI) üreterek manuel "START" ve "STOP" işlevlerini yerine getirecektir.
\end{itemize}

\subsection{Web Interface (Next.js) Requirements}

\noindent \textbf{FR-03: Live Dashboard Visualization}
\begin{itemize}
    \item Arayüz, sistemin mevcut durumunu (State) 50ms'nin altında bir gecikmeyle (latency) güncelleyerek ekranda gösterecektir.
    \item SEARCH modunda arayüzde dairesel bir radar tarama animasyonu çalışacak, TRACK modunda ise hedef kilitlenme görselleri aktifleşecektir.
\end{itemize}

\vspace{0.5cm}
\noindent \textbf{FR-04: Command \& Control Execution}
\begin{itemize}
    \item Arayüz üzerinde bulunan kontroller aracılığıyla sisteme "Sistemi Başlat", "Sistemi Durdur" ve "Manuel Hedef Kilitle" komutları gönderilebilecektir.
\end{itemize}

\vspace{0.5cm}
\noindent \textbf{FR-05: Event Logging Terminal}
\begin{itemize}
    \item Sistemden gelen tüm durum değişiklikleri ve kritik uyarılar, arayüzdeki bir terminal bileşeninde zaman damgasıyla (timestamp) listelenecektir.
\end{itemize}

\section{Communication \& Protocol Requirements (CPR)}

\noindent \textbf{CPR-01: AC2 Telemetry Protocol}\\
STM32 ve Gateway arasındaki UART haberleşmesi, veri bozulmalarını engellemek için özel bir çerçeve (Frame) formatında yapılacaktır.
\begin{itemize}
    \item \textbf{Format:} \texttt{[SYNC\_BYTE] [LENGTH] [COMMAND\_ID] [PAYLOAD] [CHECKSUM]}
    \item \textbf{SYNC\_BYTE:} Paketin başlangıcını belirten sabit byte (Örn: \texttt{0xAA}).
    \item \textbf{CHECKSUM:} Veri bütünlüğünü doğrulamak için XOR veya CRC-8 algoritması kullanılacaktır.
\end{itemize}

\vspace{0.5cm}
\noindent \textbf{CPR-02: Command Registry}\\
Desteklenen temel iletişim komutları şunlardır:
\begin{itemize}
    \item \texttt{CMD\_SET\_STATE (0x10)}: Arayüzden gömülü sisteme mod değiştirme emri.
    \item \texttt{CMD\_REPORT\_STATE (0x20)}: Gömülü sistemden arayüze anlık durum bildirimi.
    \item \texttt{CMD\_HEARTBEAT (0x99)}: Sistemin bağlantı kopmalarını tespit etmek için saniyede bir gönderdiği yaşam sinyali.
\end{itemize}

\vspace{0.5cm}
\noindent \textbf{CPR-03: WebSocket Gateway}
\begin{itemize}
    \item Bridge yazılımı, UART üzerinden aldığı AC2 Telemetry paketlerini JSON formatına ayrıştıracak ve anında WebSocket (\texttt{socket.io} vb.) üzerinden Next.js arayüzüne yayınlayacaktır.
\end{itemize}

\end{document}